\chapter{Introducción}\label{ch:introduccion}


En el panorama contemporáneo, el desarrollo de sistemas computacionales y arquitecturas orientadas a la 
gestión en el sector de la salud juega un rol determinante en la modernización de los procesos de seguimiento 
y diagnóstico médico. Específicamente dentro de la investigación dermatológica, la capacidad de almacenar, 
organizar y proveer información estructurada de manera ágil se ha convertido en una necesidad tecnológica 
indispensable. Para Cuba, el avance hacia la madurez y soberanía científico-tecnológica implica el desarrollo 
e incorporación de plataformas de \textit{software} propias, capaces de moldearse a las exigencias regulatorias, 
dinámicas de trabajo y características de su sistema de salud pública.

Surge así, como una integración interinstitucional, el Ecosistema DermaUH; un esfuerzo dirigido por la 
Facultad de Matemática y Computación de la Universidad de La Habana. Un ecosistema informático de tal 
magnitud, construido de forma multimodal para abarcar tanto interfaces web como aplicaciones desplegadas 
en terminales móviles, precisa ineludiblemente de una base centralizada. Este componente de arquitectura 
debe asumir el rol fundamental de persistir las observaciones y constituir el acervo que abastezca de 
información a las distintas capas de servicio. Disponer de este bloque fundacional no solo facilita la 
sustitución de aplicaciones importadas u onerosas, sino que además salvaguarda los datos médicos, 
considerándolo un requisito intrínseco de soberanía nacional.

Históricamente, los especialistas e investigadores han recurrido a ecosistemas foráneos y repositorios 
genéricos en un esfuerzo por suplir la ausencia de herramientas integrales en la región. Durante años, 
la investigación asistida por computadora a nivel mundial ha logrado hitos relevantes sentando principios 
de estandarización para bancos de datos dermatológicos. No obstante, al procurar 
emplear metodológicamente dichos contenedores en un entorno distinto, resulta palpable una desconexión 
metodológica y técnica.

En el país, las prácticas dermatológicas cotidianas y la investigación científica llevada a cabo por 
académicos han adolecido de una plataforma propia que cierre el ciclo de vida del dato médico: desde 
el momento exacto en que se adquiere la imagen mediante dispositivos personales en una consulta local, 
hasta su enriquecimiento a través de anotaciones de especialistas. La utilización aislada de repositorios 
internacionales\footnote{Véase el Capítulo \ref{ch:estado-del-arte} para un análisis detallado de las 
principales bases de imágenes dermatológicas globales, sus características técnicas y limitaciones.}
no satisface la necesidad de retroalimentación inmediata, ni facilita la elaboración de 
resúmenes estadísticos que consideren las divisiones demográficas de la población, un requisito vital 
para la epidemiología del país y para cimentar futuros modelos predictivos adaptados a las condiciones 
de la región. De esta desconexión entre la disponibilidad teórica de galerías de imágenes en Internet y 
la carencia empírica de un motor de bases de datos integrado a los servicios consumibles por el médico 
nacional, emerge la justificación primordial del actual desarrollo de \textit{software}.

La investigación realizada parte de una contradicción evidente: el ecosistema médico y tecnológico en Cuba carece de 
una plataforma integrada que no solo almacene imágenes clínicas y dermatoscópicas, sino que las dote de 
anotaciones de calidad, capacidades de generación de estadísticas y total integración con sistemas de 
software heterogéneos. Por lo tanto, el \textbf{problema científico} 
de esta investigación queda definido mediante la interrogante: ¿Cómo diseñar y desarrollar una base de 
imágenes dermatoscópicas anotadas que, operando como el núcleo de persistencia y análisis dentro del 
ecosistema DermaUH, garantice la fiabilidad, flexibilidad y escalabilidad necesarias para apoyar el 
diagnóstico clínico, la capacitación y la investigación biomédica en el país?


Atendiendo al problema planteado expuesto, se delimita como \textbf{objeto de estudio} el proceso de gestión, 
integración y análisis de imágenes médicas dentro de infraestructuras de \textit{software} sanitario. De manera más 
concreta, el \textbf{campo de acción} donde incide directamente este proyecto es el diseño e implementación 
de bases de datos médicas y arquitecturas de software bajo un enfoque matemático-computacional, orientados 
a la dermatología asistida.

Como \textbf{objetivo general}, se persigue desarrollar la Base de Imágenes dermatoscópicas anotadas asociada 
al ecosistema DermaUH, garantizando que conste de la flexibilidad necesaria para integrarse a aplicaciones 
externas y operar como fuente para la investigación biomédica y la visualización de estadísticas demográficas.

Para materializar este propósito, se trazan los siguientes \textbf{objetivos específicos}:
\begin{itemize}
    \item Lograr una panorámica de las principales bases de imágenes anotadas de la dermatología a 
    nivel global, examinando sus arquitecturas, servicios y esquemas de metadatos.
    \item Desarrollar y modelar una base de datos flexible, exhaustivamente anotada, que forme parte intrínseca 
    del ecosistema DermaUH, es decir, que adquiera datos y se nutra dinámicamente desde la aplicación web y la 
    aplicación móvil institucional.
    \item Lograr que este diseño arquitectónico y de datos posibilite obtener disímiles estadísticas 
    poblacionales y analíticas que sirvan como apoyo fundamental al diagnóstico experto y a la capacitación 
    de nuevos galenos.
\end{itemize}


La transformación digital del sector de la salud constituye una prioridad ineludible. En la actualidad, el 
desarrollo de historia clínica electrónica y la patología digital se posicionan como pilares en los 
sistemas sanitarios modernos~\cite{prensalatina_dermauh}. 

La novedad científica de este trabajo reside en su enfoque de ciclo cerrado, donde se abarca 
la integración multimodal de recolección \textit{in the wild} —mediante dispositivos móviles en 
consulta— hacia un repositorio estructurado que responde a esquemas formales y normalizados bajo el 
estándar del proyecto. Este paradigma se aleja del tradicional almacenamiento pasivo para concebir 
una arquitectura dinámica y orientada a servicios que permite al ecosistema no solo ver consultas y 
emitir diagnósticos en tiempo real, sino generar inteligencia acumulativa en torno a características 
demográficas específicas. 

Desde la dimensión de la importancia teórica y práctica, la tesis fortalece el \textit{corpus} metodológico 
en la ingeniería de software aplicada a la bioinformática en plataformas de escasos recursos. Su utilización 
trasciende el ejercicio académico, introduciendo un activo tecnológico que dota a la Facultad de Matemática 
y Computación y a los centros de salud asociados de un producto soberano de incalculable valor. Fomenta 
indirectamente la creación de redes neuronales sin sesgo, elevando la precisión y democratizando el acceso 
a diagnósticos preventivos de alta gama.


El transcurso de esta investigación se sustenta en la siguiente \textbf{hipótesis} de corte ingenieril y 
práctico: si se diseña e implementa una arquitectura flexible para una base de imágenes dermatoscópicas 
que integre de manera coherente metadatos clínicos locales y se acople fluidamente con las aplicaciones 
web y móviles del Ecosistema DermaUH, entonces será factible elevar la organización y trazabilidad de los 
datos dermatológicos en el país, impulsando favorablemente la precisión del diagnóstico clínico, el 
suministro de datos estadísticos fiables y la capacitación efectiva de especialistas de la salud.

Para un mejor abordaje y comprensión del proyecto, el presente documento se encuentra divido en secciones 
claramente definidas. El \textbf{Capítulo \ref{ch:estado-del-arte}} aborda el Estado del Arte, estudiando 
sistémicamente los esfuerzos mundiales en repositorios dermatológicos y estableciendo paralelos técnicos. 
Seguidamente, en el \textbf{Capítulo \ref{ch:proposal}} se expone la Propuesta, abarcando el sustento lógico, 
los flujogramas de información, las funcionalidades con las que cuenta el sistema, los esquemas de bases 
de datos y la estructuración del diseño 
teórico-práctico. El \textbf{Capítulo \ref{ch:implementation}} detalla la Implementación tecnológica, 
precisando las herramientas, los lenguajes, los resultados algorítmicos tangibles en el contexto del 
ecosistema y las pruebas funcionales realizadas. Finalmente, la obra cierra enunciando de forma sintetizada las conclusiones a las que se arriban 
tras el desarrollo funcional de la herramienta, e incluye recomendaciones prospectivas que delinean la hoja de 
ruta de la investigación en sus siguientes etapas funcionales. Además, se agrega un apartado con la 
bibliografía y las referencias bibliográficas utilizadas como sustento del trabajo.

Finalmente, cabe destacar que a lo largo de la presente investigación se incorporó el uso de herramientas 
de Inteligencia Artificial Generativa, tanto en la redacción del informe escrito como en el desarrollo 
del producto de \textit{software}. En ambos escenarios, estas tecnologías fueron empleadas fundamentalmente 
para concebir la estructura base inicial deseada, proporcionando un punto de partida organizativo sólido. 
En lo concerniente al sistema informático, su utilización se enfocó de manera mayoritaria en el diseño visual de 
la aplicación y en la construcción de su interfaz gráfica. De igual modo, se recurrió a estas herramientas 
como recurso de apoyo para la corrección y depuración de errores de baja complejidad que ya habían sido 
previamente identificados, contribuyendo a la optimización del proceso de desarrollo sin comprometer 
el rigor técnico y metodológico del ecosistema.
